\chapter{Technischer Ist-Zustand}\label{ch:technischer-ist-zustand}

\section{Vorbereitung}\label{sec:vorbereitung}
Zur Vorbereitung der Erfassung des Technischen-ist-zustandes wurden überlegungen angestellt, wie sich der Prozess des Zuordnens einer ID, sowie relevanten Informationen sowie einem Bild, falls nötig, möglichst strukturiert und übersichtlich darstelen lässt.
Ebenfalls relevent hier war es, dass die Daten möglichst leicht in eine \LaTeX-Tabelle überführt werden können.
Hier wurden drei mögliche Konzepte, die zur effizienten aufnahme der Daten in Betracht gezogen wurden.

Für die Excel-Datei, die Web-App und die APK wurde Generative KI eingesetzt.
Diese hat lediglich ein Großteil des Programmcodes, sowie das Hosting der Web-App übernommen.
Ebenfalls wurden Formulierungen angepasst und Texte auf Gramatik und Rechtschreibung überprüft.
Alle bewertenden Elemente, sowie die Begründung wurde durch ein Gruppenmitglied erstellt und nach überarbeitung und umformulierung auf den vorgegebenen Inhalt überprüft.
Alle Entscheidungen wurden ebenfalls aufgrund der eigenen Bewertungsmatrix getroffen.

Die Excel-Datei bietet eine einfache Einrichtung, eine zuverlässige Offline-Nutzung und einen unkomplizierten Datenexport.
Die Bedienung auf einem Tablet ist aufgrund des Touchsystems allerdings sehr unzuverlässig, da man sich schnell mal vertippt, bzw was Falsches auswählen kann.
Außerdem lassen sich komplexe Verbindungen zwischen Bauteilen und zugehörige Fotos nur eingeschränkt übersichtlich erfassen.
Hinzu kommt, dass es für Mobile Endgeräte nur die Office 365 Version von Excel gibt, welche auf Android geräten in ihren Funktionen stark beschränkt ist

Die Web-App ist auf die konkrete Erfassungsaufgabe zugeschnitten und bietet automatische IDs, getrennte Ansichten für Bauteile und Verbindungen sowie große, für Touchscreens geeignete Eingabeelemente.
Ihre Schwächen liegen in der Abhängigkeit vom Browser, der lokalen Speicherung und dem nicht auf jedem Gerät gleichermaßen zuverlässigen Offlinebetrieb.
Ebenfalls zu beachten ist hier, dass man jene erst erstellen muss, welches ein gewisses Grundverständnis von Programmierung mit gezielter nutzung von KI erfordert.

Die APK übernimmt die aufgabenbezogene Bedienoberfläche der Web-App und stellt sie als eigenständige Android-Anwendung bereit.
Sie ermöglicht eine zuverlässige Offline-Nutzung und einen direkten Zugriff auf Gerätefunktionen.
Außerdem kann so eine Umgebung geschaffen werden, die die exakten Bedingungen für das hier geplante vorgehen bereitstellt.
Dem stehen ein höherer Aufwand bei der erstmaligen Erstellung gegenüber.
Aufgrund der hohen Bedienbarkeit und der guten Eignung für die mobile Datenerfassung wurde die APK als bevorzugte Lösung ausgewählt.

Für einen übersichtlichen vergleich werden die Optionen in \cref{tab:nutzwertanalyse-erfassung} aufgelistet und gegeneinnander abgewogen und dann anschließend mit der Formel
\begin{equation}
    N_j =
    \frac{\sum_{i=1}^{n} w_i \cdot b_{ij}}
    {b_{\max} \cdot \sum_{i=1}^{n} w_i}
    \cdot 100\,\%\label{eq:equation}
\end{equation}
bestimmt werden, wobei $w_i$ hier die Prozentuale Gewichtung ist und $b_{ij}$ die selbstgewählte Bewertung, welche von $1 - 5$ gesetzt wurde.

\begin{table}[htbp]
    \centering
    \caption{Nutzwertanalyse der betrachteten Erfassungslösungen}
    \label{tab:nutzwertanalyse-erfassung}
    \begin{tabular}{p{5.3cm}rrrr}
        \hline
        \textbf{Kriterium} &
        \textbf{Gewichtung} &
        \textbf{Excel} &
        \textbf{Web-App} &
        \textbf{APK} \\
        \hline

        Bedienbarkeit auf dem Tablet
        & 20\,\% & 2 & 4 & 5 \\

        Offline-Nutzung
        & 15\,\% & 5 & 3 & 5 \\

        Automatische IDs und Fehlervermeidung
        & 15\,\% & 3 & 5 & 5 \\

        Erfassung von Verbindungen
        & 10\,\% & 2 & 5 & 5 \\

        Einbindung von Fotos
        & 10\,\% & 1 & 3 & 5 \\

        Export für \LaTeX
        & 10\,\% & 5 & 5 & 5 \\

        Einrichtung und Installation
        & 10\,\% & 3 & 4 & 4 \\

        Sicherung und Datenübertragung
        & 5\,\% & 5 & 3 & 4 \\

        Änderungen und Aktualisierungen
        & 5\,\% & 4 & 1 & 1 \\
        \hline

        \textbf{Gesamtbewertung}
        & \textbf{100\,\%}
        & \textbf{63\,\%}
        & \textbf{78\,\%}
        & \textbf{93\,\%} \\

        \hline
    \end{tabular}
\end{table}

mit
\begin{align}
    N_{\text{Excel}}
    &=
    \frac{
        20 \cdot 2 +
        15 \cdot 5 +
        15 \cdot 3 +
        10 \cdot 2 +
        10 \cdot 1 +
        10 \cdot 5 +
        10 \cdot 3 +
        5 \cdot 5 +
        5 \cdot 4
    }{500}
    \cdot 100\,\% \\
    &=
    \frac{315}{500}
    \cdot 100\,\%
    =
    63\,\%,
    \\[1em]
    N_{\text{Web-App}}
    &=
    \frac{
        20 \cdot 4 +
        15 \cdot 3 +
        15 \cdot 5 +
        10 \cdot 5 +
        10 \cdot 3 +
        10 \cdot 5 +
        10 \cdot 4 +
        5 \cdot 3 +
        5 \cdot 1
    }{500}
    \cdot 100\,\% \\
    &=
    \frac{390}{500}
    \cdot 100\,\%
    =
    78\,\%,
    \\[1em]
    N_{\text{APK}}
    &=
    \frac{
        20 \cdot 5 +
        15 \cdot 5 +
        15 \cdot 5 +
        10 \cdot 5 +
        10 \cdot 5 +
        10 \cdot 5 +
        10 \cdot 4 +
        5 \cdot 4 +
        5 \cdot 1
    }{500}
    \cdot 100\,\% \\
    &=
    \frac{465}{500}
    \cdot 100\,\%
    =
    93\,\%.
\end{align}


Im Folgenden werden die Bewertungen der betrachteten Lösungen anhand der zuvor festgelegten Kriterien begründet.


Die Bedienbarkeit erhielt die höchste Gewichtung, da sie einen unmittelbaren Einfluss auf die Geschwindigkeit und Fehleranfälligkeit der Datenerfassung hat.
Bei einer großen Anzahl von Bauteilen können umständliche oder häufig wiederkehrende Eingabeschritte den Arbeitsablauf erheblich verlangsamen.

Die Offline-Nutzung und der Automatisierungsgrad wurden ebenfalls als wichtige Kriterien berücksichtigt.
Eine zuverlässige Offline-Funktion ermöglicht die Erfassung unabhängig von der Qualität der verfügbaren Internetverbindung und verringert das Risiko, dass Arbeitsfortschritte aufgrund von Verbindungsproblemen nicht vollständig gespeichert werden.
Da die Aufnahme einer großen Anzahl von Bauteilen viele gleichartige Arbeitsschritte umfasst, können zudem Unaufmerksamkeit und Ermüdung die Fehlerwahrscheinlichkeit erhöhen.
Eine automatische und fortlaufende Vergabe der Identifikationsnummern sowie vorgegebene Eingabe- und Auswahlmöglichkeiten reduzieren dieses Fehlerpotenzial und gewährleisten eine konsistente Datenerfassung.

Die Unterstützung bei der Erfassung von Verbindungen, die Einbindung von Fotos, der Export für \LaTeX{} sowie der Einrichtungsaufwand wurden im Vergleich zu Bedienbarkeit, Offline-Nutzung und Automatisierung geringer gewichtet.
Diese Funktionen erleichtern den Arbeitsablauf, sind für die grundlegende Aufnahme der Bauteile jedoch nicht in jedem Fall zwingend erforderlich und könnten bei Bedarf durch manuelle Arbeitsschritte ersetzt werden.
Da die Erfassung einer großen Anzahl von Elementen viele wiederkehrende Tätigkeiten umfasst, bieten solche Komfortfunktionen dennoch einen deutlichen Mehrwert.
Sie reduzieren den erforderlichen Arbeitsaufwand, vereinfachen die Dokumentation und senken die Wahrscheinlichkeit von Fehlern, die durch Unaufmerksamkeit oder Ermüdung entstehen können.

Die Kriterien „Sicherung und Datenübertragung“ sowie „Änderungen und Aktualisierungen“ erhielten eine vergleichsweise geringe Gewichtung.
Der Export und die Übertragung der erfassten Daten stellen bei allen betrachteten Lösungen grundsätzlich keinen größeren Aufwand dar.
Eine regelmäßige Sicherung ist zwar wichtig, beeinflusst den eigentlichen Vorgang der Katalogisierung jedoch nur indirekt.
Auch die nachträgliche Anpassbarkeit wurde geringer gewichtet, da die Excel-Datei, die Web-App und die Android-App bereits vor ihrem praktischen Einsatz getestet und auf mögliche Verbesserungen überprüft werden.
Dennoch können während der Erfassung unerwartete Probleme oder neue Anforderungen auftreten.
In diesem Fall unterscheiden sich die Lösungen darin, wie schnell Änderungen vorgenommen und vor Ort bereitgestellt werden können.
Insbesondere bei der Android-App ist dafür die Erstellung und Installation einer aktualisierten APK erforderlich.

Bei der Bewertung von Excel ist außerdem zu berücksichtigen, dass es sich um eine herstellerspezifische Software des Unternehmens Microsoft handelt.
Der Quellcode ist nicht öffentlich zugänglich, und die Weiterentwicklung, Lizenzierung sowie langfristige Verfügbarkeit der Anwendung werden durch den Hersteller bestimmt.
Je nach eingesetzter Version können zudem ein Microsoft-Konto, eine kostenpflichtige Lizenz oder die Nutzung von Cloud-Diensten erforderlich sein.
Daraus entstehen mögliche Abhängigkeiten vom Anbieter und Fragestellungen hinsichtlich Datenschutz und Datenhoheit.
Da die Tabellen grundsätzlich lokal gespeichert und in offene Formate wie CSV exportiert werden können, wurde Excel dennoch als mögliche Lösung berücksichtigt.

\newpage
\subsubsection{Excel-Datei}

\textbf{Bedienbarkeit auf dem Tablet (2/5):}
\begin{itemize}[leftmargin=2em, nosep]
    \item[\pluspunkt] Die tabellarische Darstellung bietet eine kompakte
    Übersicht über die erfassten Daten.
    \item[\minuspunkt] Kleine Zellen und Bedienelemente erschweren die
    Bedienung auf einem Touchscreen.
    \item[\minuspunkt] Wiederkehrende Eingaben und der Wechsel zwischen
    Zellen verlangsamen den Arbeitsablauf.
    \item[\minuspunkt]
    Schäbiges App Layout.
\end{itemize}

\medskip
\textbf{Offline-Nutzung (5/5):}
\begin{itemize}[leftmargin=2em, nosep]
    \item[\pluspunkt] Eine lokal gespeicherte Datei kann vollständig ohne
    Internetverbindung bearbeitet werden.
\end{itemize}

\medskip
\textbf{Automatische IDs und Fehlervermeidung (3/5):}
\begin{itemize}[leftmargin=2em, nosep]
    \item[\pluspunkt] IDs können durch Formeln oder vorbereitete
    Zahlenreihen erzeugt werden.
    \item[\minuspunkt] Formeln und IDs können versehentlich verändert,
    gelöscht oder überschrieben werden.
    \item[\minuspunkt]
    In der Excel App sind manche Formeln und Funktionen nicht möglich, was das Automatisieren beschränkt.
\end{itemize}

\medskip
\textbf{Erfassung von Verbindungen (2/5):}
\begin{itemize}[leftmargin=2em, nosep]
    \item[\pluspunkt] Verbindungen können in einem zweiten Tabellenblatt
    dokumentiert werden.
    \item[\minuspunkt] Die Auswahl und Zuordnung der beteiligten Bauteile
    ist auf dem Tablet vergleichsweise umständlich.
    \item[\minuspunkt]
    Man muss zwischen Tabs wechseln, was den Workflow stört.
    \item[\minuspunkt]
    Die IDs sind nicht vorbereitet, man muss alle selber eingeben.
\end{itemize}

\medskip
\textbf{Einbindung von Fotos (1/5):}
\begin{itemize}[leftmargin=2em, nosep]
    \item[\minuspunkt] Fotos lassen sich nur umständlich in Tabellenzellen
    einfügen und den Bauteilen zuordnen.
    \item[\minuspunkt] Eingefügte Bilder können die Dateigröße und
    Übersichtlichkeit deutlich beeinträchtigen.
    \item[\minuspunkt]
    Das Bewegen von Fotos ist aufgrund der Touchfunktion und die automatische skalierung von Excel ebenfalls komplizierter.
\end{itemize}


\medskip
\textbf{Export für \LaTeX{} (5/5):}
\begin{itemize}[leftmargin=2em, nosep]
    \item[\pluspunkt] Tabellenblätter können direkt als CSV-Dateien
    exportiert und in \LaTeX{} eingelesen werden.
\end{itemize}

\medskip
\textbf{Einrichtung und Installation (3/5):}
\begin{itemize}[leftmargin=2em, nosep]
    \item[\pluspunkt] Nach Erstellung der Tabelle kann diese theoretisch unmittelbar
    verwendet werden.
    \item[\minuspunkt] Auswahlfelder, Formatierungen und Prüfregeln müssen
    zunächst eingerichtet werden.
    \item[\minuspunkt] Excel ist ein herstellerspezifische Produkt von Microsoft, welches einen Account benötigt, der aufgrund eines Problems mit der Studi Lizenz nur erschwert zu bekommen ist.
\end{itemize}

\medskip
\textbf{Sicherung und Datenübertragung (5/5):}
\begin{itemize}[leftmargin=2em, nosep]
    \item[\pluspunkt] Die gesamte Erfassung befindet sich in einer einzelnen
    Datei, die einfach kopiert und gesichert werden kann.
\end{itemize}

\medskip
\textbf{Änderungen und Aktualisierungen (4/5):}
\begin{itemize}[leftmargin=2em, nosep]
    \item[\pluspunkt] Spalten, Kategorien und Auswahlfelder können
    nachträglich angepasst werden.
    \item[\minuspunkt] Änderungen an der Tabellenstruktur können bestehende
    Formeln oder Verweise beeinträchtigen.
\end{itemize}

\newpage

\subsubsection{Web-App}

\textbf{Bedienbarkeit auf dem Tablet (4/5):}
\begin{itemize}[leftmargin=2em, nosep]
    \item[\pluspunkt] Große Eingabefelder und Schaltflächen erleichtern die
    Bedienung auf einem Touchscreen.
    \item[\pluspunkt] Bauteile und Verbindungen werden in getrennten
    Ansichten dargestellt.
    \item[\minuspunkt] Die Bedienung erfolgt innerhalb eines Browsers und
    kann durch dessen Benutzeroberfläche beeinflusst werden.
\end{itemize}

\medskip
\textbf{Offline-Nutzung (3/5):}
\begin{itemize}[leftmargin=2em, nosep]
    \item[\pluspunkt] Nach der ersten Einrichtung kann die Anwendung
    grundsätzlich offline verfügbar gemacht werden.
    \item[\minuspunkt] Der Offlinebetrieb hängt vom Browser und dessen
    Zwischenspeicher ab.
\end{itemize}

\medskip
\textbf{Automatische IDs und Fehlervermeidung (5/5):}
\begin{itemize}[leftmargin=2em, nosep]
    \item[\pluspunkt] Eindeutige IDs werden automatisch vergeben und bleiben
    beim Bearbeiten eines Bauteils erhalten.
    \item[\pluspunkt] Auswahlfelder verhindern viele fehlerhafte oder
    uneinheitliche Eingaben.
\end{itemize}

\medskip
\textbf{Erfassung von Verbindungen (5/5):}
\begin{itemize}[leftmargin=2em, nosep]
    \item[\pluspunkt] Bereits erfasste Bauteile können anhand ihrer ID und
    Bezeichnung ausgewählt werden.
    \item[\pluspunkt] Die Anwendung verhindert die Verbindung eines Bauteils
    mit sich selbst.
    \item[\pluspunkt] Die Anwendung bietet bei Verbindungen nur bislang erstellte IDs an.
\end{itemize}

\medskip
\textbf{Einbindung von Fotos (3/5):}
\begin{itemize}[leftmargin=2em, nosep]
    \item[\pluspunkt] Fotos können direkt ausgewählt und einem Bauteil
    zugeordnet werden.
    \item[\minuspunkt] Die Speicherung vieler Fotos kann den verfügbaren
    Browserspeicher stark beanspruchen.
    \item[\minuspunkt] Die UI für die Fotoauswahl im Browser ist nicht top of the shelf.
\end{itemize}

\medskip
\textbf{Export für \LaTeX{} (5/5):}
\begin{itemize}[leftmargin=2em, nosep]
    \item[\pluspunkt] Bauteile und Verbindungen können getrennt als
    CSV-Dateien exportiert werden.
    \item[\pluspunkt] Es kann ein Backup erstellt werden, was das arbeiten auch auf mehrere Tage problemslos aufteilen lässt .
\end{itemize}

\medskip
\textbf{Einrichtung und Installation (4/5):}
\begin{itemize}[leftmargin=2em, nosep]
    \item[\pluspunkt] Die Anwendung Anwendung kann über einen Link geöffnet
    werden und benötigt keine klassische Installation.
    \item[\minuspunkt] Für den ersten Zugriff und die Vorbereitung des
    Offlinebetriebs ist eine Internetverbindung erforderlich.
    \item[\minuspunkt] Es ist ein OpenAI Account notwendig, um auf die Wesite zuzugreifen.
\end{itemize}

\medskip
\textbf{Sicherung und Datenübertragung (3/5):}
\begin{itemize}[leftmargin=2em, nosep]
    \item[\pluspunkt] Die erfassten Daten können als Sicherungsdatei
    heruntergeladen werden.
    \item[\minuspunkt] Die Daten sind zunächst an den jeweiligen Browser
    gebunden und müssen manuell gesichert werden.
\end{itemize}

\medskip
\textbf{Änderungen und Aktualisierungen (1/5):}
\begin{itemize}[leftmargin=2em, nosep]
    \item[\minuspunkt] Die dokumentierende Person kann Aufbau und Funktionen
    der Anwendung nicht selbstständig verändern.
    \item[\minuspunkt] Anpassungen erfordern Änderungen am Programmcode und
    eine erneute Veröffentlichung.
\end{itemize}

\newpage

\subsubsection{Android-APK}

\textbf{Bedienbarkeit auf dem Tablet (5/5):}
\begin{itemize}[leftmargin=2em, nosep]
    \item[\pluspunkt] Die Oberfläche wurde gezielt für die Bedienung auf
    einem Android-Tablet entwickelt.
    \item[\pluspunkt] Große Eingabefelder, Schaltflächen und getrennte
    Ansichten ermöglichen einen schnellen Arbeitsablauf.
\end{itemize}

\medskip
\textbf{Offline-Nutzung (5/5):}
\begin{itemize}[leftmargin=2em, nosep]
    \item[\pluspunkt] Alle benötigten Anwendungsdateien befinden sich
    innerhalb der APK und stehen ohne Internetverbindung zur Verfügung.
\end{itemize}

\medskip
\textbf{Automatische IDs und Fehlervermeidung (5/5):}
\begin{itemize}[leftmargin=2em, nosep]
    \item[\pluspunkt] Die IDs werden automatisch und eindeutig vergeben.
    \item[\pluspunkt] Auswahlfelder und Eingabeprüfungen reduzieren
    fehlerhafte Einträge.
\end{itemize}

\medskip
\textbf{Erfassung von Verbindungen (5/5):}
\begin{itemize}[leftmargin=2em, nosep]
    \item[\pluspunkt] Bauteile können bequem über ihre ID und Bezeichnung
    ausgewählt und miteinander verknüpft werden.
    \item[\pluspunkt] Die Verbindungen werden getrennt von den Bauteildaten
    gespeichert und übersichtlich dargestellt.
\end{itemize}

\medskip
\textbf{Einbindung von Fotos (5/5):}
\begin{itemize}[leftmargin=2em, nosep]
    \item[\pluspunkt] Fotos können direkt über die Funktionen des Tablets
    ausgewählt und einem Bauteil zugeordnet werden.
    \item[\pluspunkt] Die aufgenommenen Bilder werden gemeinsam mit den
    übrigen Daten gesichert.
    \item[\pluspunkt] UI zum Auswählen ist sehr nutzerfreundlich.
\end{itemize}

\medskip
\textbf{Export für \LaTeX{} (5/5):}
\begin{itemize}[leftmargin=2em, nosep]
    \item[\pluspunkt] Bauteile und Verbindungen können als getrennte
    CSV-Dateien exportiert werden.
    \item[\pluspunkt] Der Android-Dialog ermöglicht das Speichern und Teilen
    der erzeugten Dateien.
\end{itemize}

\medskip
\textbf{Einrichtung und Installation (4/5):}
\begin{itemize}[leftmargin=2em, nosep]
    \item[\pluspunkt] Nach einmaliger Installation kann die Anwendung direkt
    über das App-Symbol gestartet werden.
    \item[\minuspunkt] Da die Anwendung nicht über den Play Store verteilt
    wird, muss die Installation aus einer externen Quelle einmalig erlaubt
    werden.
\end{itemize}

\medskip
\textbf{Sicherung und Datenübertragung (4/5):}
\begin{itemize}[leftmargin=2em, nosep]
    \item[\pluspunkt] Sicherungsdateien können über den Android-Dialog
    gespeichert oder an ein anderes Gerät übertragen werden.
    \item[\minuspunkt] Die Sicherung muss durch die dokumentierende Person
    manuell durchgeführt werden.
\end{itemize}

\medskip
\textbf{Änderungen und Aktualisierungen (1/5):}
\begin{itemize}[leftmargin=2em, nosep]
    \item[\minuspunkt] Änderungen erfordern eine Anpassung des Programmcodes
    und die Erstellung einer neuen APK, was lokal nicht möglich ist.
\end{itemize}